\section{Code Optimization}
\label{sec:pie}

Performance-Improving Code Edits or PIE~\citep{pie} focuses on enhancing the efficiency of functionally correct programs.
The primary objective of PIE is to optimize a given program by implementing algorithmic modifications that lead to improved runtime performance.

Given an optimization generated by PIE, \ours first generates a natural language feedback on possible improvements~\Cref{fig:prompt:pie:fb}.
Then, the feedback is fed to \itermod~\Cref{fig:prompt:pie:refine} for refinement.


\begin{table}[ht]
    \centering
    \caption{Main Results and Ablation Analysis}
    \begin{tabular}{lccccc}
        \toprule
        Setup          & Iteration & \% Optimized & Relative Speedup & Speedup \\
        \midrule
        Direct         & -         & 9.7          & 62.29           & 3.09    \\
        \midrule
        \ours $-$ feedback & 1         & 10.1         & 62.15           & 3.03    \\
        \ours $-$ feedback & 2         & 10.4         & 61.79           & 3.01    \\
        \midrule
        \ours   & 1         & 15.3         & 59.64           & 2.90    \\
        \ours   & 2         & \textbf{15.6}         & \textbf{65.60 }          & \textbf{3.74}    \\
        \bottomrule
    \end{tabular}
    \caption{Performance comparison of \ours and ablated variants for code optimization. The table highlights the effectiveness of \ours in optimizing code through iterative feedback and improvement, outperforming both the direct method and the simplified feedback approach, which lacks the introspective feedback mechanism of \ours. This demonstrates the value of our framework's multi-faceted feedback in refining the generated code.}
    \label{tab:results_ablation}
\end{table}


